\documentclass[12pt,a4paper]{article}

% ==== Paquetes básicos ====
\usepackage[spanish,es-noquoting,es-tabla]{babel}
\usepackage[utf8]{inputenc}
\usepackage[T1]{fontenc}
\usepackage{listings}
\usepackage{fvextra}      % mejora UTF-8 en listings y fija bugs internos
\usepackage{xcolor}
\usepackage{lmodern}
\usepackage{geometry}
\usepackage{setspace}
\usepackage{microtype}
\usepackage{siunitx}
\usepackage{graphicx}
\usepackage{subcaption}
\usepackage{array}
\usepackage{booktabs}
\usepackage{multirow}
\usepackage{caption}
\usepackage{hyperref}
\usepackage{enumitem}
\usepackage{titlesec}
\usepackage{amsmath,amssymb}
\usepackage{fancyhdr}
\usepackage{float}
\usepackage[section]{placeins}

% ==== Configuración de página y tipografía ====
\geometry{margin=2.5cm}
\onehalfspacing
\hypersetup{
  colorlinks=true,
  linkcolor=blue,
  citecolor=blue,
  urlcolor=blue,
  pdfauthor={Antonio Enrique Coll Meseguer},
  pdftitle={Práctica de Modelado Molecular - Tripéptido Ala–Arg–Asn (ARN)},
  pdfsubject={Dinámica Molecular a 298 K y 400 K con GROMACS},
  pdfcreator={LaTeX}
}
\captionsetup{labelfont=bf}
\titlespacing*{\section}{0pt}{*3}{*1}
\titlespacing*{\subsection}{0pt}{*2}{*0.8}
\titlespacing*{\subsubsection}{0pt}{*1.6}{*0.6}
\setlength{\parskip}{0.5em}

% ==== Encabezado/pie de página sencillo ====
\pagestyle{fancy}
\setlength{\headheight}{30pt}
\fancyhf{}
\lhead{UNIVERSIDAD DE MURCIA — FACULTAD DE BIOLOGÍA — MÁSTER DE BIOINFORMÁTICA}
\rhead{Práctica de Modelado Molecular}
\cfoot{\thepage}

% ==== Listados de comandos/código ====
\lstdefinestyle{bashstyle}{
  language=bash,
  basicstyle=\ttfamily\small,
  keywordstyle=\bfseries\color{blue!70!black},
  commentstyle=\itshape\color{green!40!black},
  stringstyle=\color{red!60!black},
  showstringspaces=false,
  frame=single,
  rulecolor=\color{black!20},
  backgroundcolor=\color{black!3},
  breaklines=true,
  columns=fullflexible,
  keepspaces=true,         % respeta espacios en blanco
  numbers=left,            % (opcional) numera líneas
  numberstyle=\tiny\color{black!50},
  stepnumber=1,
  numbersep=6pt,
}


\lstset{
    style=bashstyle,
    texcl=false,         % ❗ NO enviar comentarios a LaTeX -> # y _ funcionan SIEMPRE
    mathescape=false,    % ❗ evita modo matemático dentro del bash
    inputencoding=utf8,  % UTF8 real
    extendedchars=true,
    upquote=true,        % comillas correctas
    literate=%
      {á}{{\'a}}1 {é}{{\'e}}1 {í}{{\'i}}1 {ó}{{\'o}}1 {ú}{{\'u}}1
      {Á}{{\'A}}1 {É}{{\'E}}1 {Í}{{\'I}}1 {Ó}{{\'O}}1 {Ú}{{\'U}}1
      {ñ}{{\~n}}1 {Ñ}{{\~N}}1
      {…}{{...}}1 {–}{{-}}1 {—}{{-}}1
      {¡}{{!}}1 {¿}{{?}}1
}

% ==== Portada ====
\begin{document}

\begin{titlepage}
\centering

% --- Encabezado institucional ---
{\fontsize{20}{22}\selectfont \textbf{UNIVERSIDAD DE MURCIA}}\\[4pt]
{\Large FACULTAD DE BIOLOGÍA}\\[2pt]
{\large MÁSTER DE BIOINFORMÁTICA}\\[0.8cm]

% --- LOGO o IMAGEN ---
\includegraphics[width=0.80\textwidth]{img/SelloUMU-positivo.png}\\[1 cm]

% --- Título principal ---
{\fontsize{26}{30}\selectfont \textbf{Práctica de Modelado Molecular}}\\[10pt]
{\Large Tripéptido ARN (ACE–Ala–Arg–Asn–NME)}\\[1.5cm]

% --- Datos del informe ---
\begin{tabular}{rl}
\textbf{Asignatura:} & Modelado Molecular \\
\textbf{Fecha:} & 22 de marzo de 2026 \\
\textbf{Autor:} & Antonio Enrique Coll Meseguer \\
\textbf{Curso:} & 2025/2026 \\
\end{tabular}

\vfill
\end{titlepage}

\begin{abstract}
\textbf{Resumen}. Se analiza el tripéptido ARN (ACE–Ala–Arg–Asn–NME) mediante dinámica molecular (MD) con GROMACS a \SI{298}{K} y \SI{400}{K}. El flujo de trabajo se divide en tres bloques exhaustivamente comentados: (i) \textit{EQUILIBRIO} (preparación, neutralización y equilibración NVT), (ii) \textit{RUN} (producción NVT con continuidad física) y (iii) \textit{ANÁLISIS} (Rg, distancias, ángulos, diedros, H-bonds, energía, temperatura, velocidades). Se explican todos los comandos usados, su significado y parámetros, y se proponen pies de figura para cada resultado.
\end{abstract}

\tableofcontents
\newpage

% =============================
% 1. Introducción
% =============================
\section{Introducción}
Las simulaciones se llevaron a cabo utilizando \textbf{GROMACS} (\textit{GROningen MAchine for Chemical Simulations}), un paquete empleado para dinámica molecular biomolecular. Se empleó el campo de fuerza \textbf{CHARMM27}, disponible en \texttt{share/gromacs/top/charmm27.ff}, que proporciona parámetros validados para péptidos y proteínas.

La dinámica molecular (MD) es una herramienta para estudiar biomoléculas, pues permite describir su evolución temporal a escala atómica integrando las ecuaciones de Newton. A diferencia de técnicas experimentales que proporcionan medidas promediadas, MD ofrece acceso directo a fluctuaciones conformacionales, estabilidad estructural, flexibilidad interna y efectos termodinámicos bajo condiciones controladas de temperatura y composición.

En este trabajo se evalúa el comportamiento del tripéptido ARN (ACE-Ala-Arg-Asn-NME), un modelo sencillo que permite un muestreo computacional eficiente, pero que conserva elementos estructurales esenciales del esqueleto peptídico (enlaces covalentes, ángulos, torsiones y puentes de hidrógeno).

En este contexto, se emplea MD para analizar cómo la temperatura modula la conformación, la estabilidad y la dinámica interna del tripéptido ARN en solución acuosa.

\subsubsection*{Repositorio del proyecto}
El código fuente utilizado en el postprocesado (scripts de Python), así como todas
las figuras generadas, los notebooks, los archivos auxiliares y la estructura del
proyecto completo, están disponibles públicamente en GitHub:
\url{https://github.com/antonioenriquecoll-design/modelo-molecular-ARN}.

% =============================
% 2. Objetivos
% =============================
\section{Objetivos}
El objetivo general de este trabajo es realizar y analizar simulaciones de dinámica molecular en ensamble NVT para el tripéptido ARN (ACE–Ala–Arg–Asn–NME) en disolución acuosa, siguiendo estrictamente los procedimientos desarrollados durante las sesiones prácticas de Modelado Molecular.

De forma específica, se pretende:

\begin{itemize}
  \item Ejecutar simulaciones a \SI{298}{K} y \SI{400}{K} reproduciendo los pasos aplicados en clase: generación de la topología (ACE–R1–R2–R3–NME), definición de la caja, solvatación, neutralización, equilibración NVT y producción NVT.
  
  \item Representar y comentar la evolución temporal de las siguientes magnitudes:
    \begin{itemize}
      \item energía total y energía cinética,
      \item temperatura instantánea del sistema,
      \item dos distancias internas entre átomos contiguos del esqueleto,
      \item dos ángulos internos,
      \item dos diedros representativos del backbone,
      \item radio de giro,
      \item velocidades atómicas de cinco átomos seleccionados del péptido.
    \end{itemize}

  \item Comparar cuantitativamente los resultados obtenidos a ambas temperaturas y discutir su significado estructural, dinámico y termodinámico.

  \item Evaluar la energía cinética media del sistema a \SI{298}{K} y \SI{400}{K} y contrastarla con el valor esperado según el teorema de equipartición.

  \item Garantizar que las unidades, escalas, formatos gráficos y leyendas sigan el estilo propio de artículos científicos, y comentar de forma rigurosa tanto los comandos utilizados como la información mostrada en cada gráfica.

\end{itemize}


% =============================
% 2. Preparación del sistema
% =============================
\section{Preparación del sistema}
\label{sec:metodologia}

\subsection{Preparación inicial y generación de la topología}
\textbf{Contexto}. El proyecto se llevó a cabo en el directorio \texttt{~/ARN\_project}. Todos los pasos iniciales se realizan en la raíz del proyecto y, posteriormente, se crean las carpetas de 298 K y 400 K.

\subsubsection*{Generación de la topología molecular (\texttt{gmx pdb2gmx})}
\noindent\textbf{Nota.} Durante la ejecución de \texttt{gmx pdb2gmx -ter}, se seleccionó el campo de fuerza \textbf{CHARMM27} (número 8)y el modelo de agua \textbf{TIP3P} (número 1). Para los terminales del tripéptido se escogió la opción \texttt{None} tanto en el extremo N (número 2) como en el C (número 4), de forma que se respetan los grupos terminales definidos en el PDB. Estas elecciones generan la topología estándar compatible con CHARMM27.
\begin{lstlisting}[style=bashstyle,texcl=false]
# Crear y entrar en la carpeta raiz del proyecto
mkdir -p ~/ARN_project
cd ~/ARN_project

# Copiar el PDB inicial del tripeptido ARN
cp ../for-students/tripeptides/arn/arn.pdb .

# Generar coordenadas (.gro) y topologia (.top) con seleccion manual de terminales
# -f: entrada PDB; -o: salida .gro; -p: salida .top; -ter: selección de terminales
# Campo de fuerza: CHARMM27; Agua: TIP3P
gmx pdb2gmx -f arn.pdb -o arn.gro -p arn.top -ter
\end{lstlisting}

\textbf{Qué hace}. Asigna tipos atómicos, cargas y conectividades compatibles
con CHARMM27, produciendo \texttt{arn.gro} y \texttt{arn.top}. \textbf{Salida}.
Reconoce ACE/ALA/ARG/ASN/NME; la topología incluye secciones~\verb+[ atoms ]+,
enlaces, ángulos y dihedros.

\subsubsection*{Inspección inicial del archivo \texttt{arn.pdb} mediante \texttt{less}}

Antes de generar la topología con \texttt{pdb2gmx}, se inspeccionó el archivo
\texttt{arn.pdb} para identificar la organización estructural del tripéptido
original en formato PDB:

\begin{verbatim}
less arn.pdb
\end{verbatim}

Un fragmento relevante del fichero muestra:

\begin{verbatim}
ATOM      7   N   ALA     2      1.662   1.368  -0.150
ATOM      8   CA  ALA     2      3.068   1.444  -0.538
ATOM      9   C   ALA     2      3.850   2.297   0.431
ATOM     17   N   ARG     3      4.995   1.877   0.851
ATOM     18   CA  ARG     3      5.752   2.702   1.788
\end{verbatim}

El formato de cada línea del PDB es:

\begin{center}
\texttt{ATOM  Index  AtomName  Residue  ResID   x   y   z}
\end{center}

Donde las coordenadas están en Ångström.

\paragraph*{Importante.}
Los índices del PDB \textbf{no coinciden} con los índices finales que utiliza GROMACS,
ya que el programa \texttt{pdb2gmx}:
\begin{itemize}
  \item añade hidrógenos explícitos,
  \item reordena los átomos según el campo de fuerza,
  \item genera nuevos números de átomo para el archivo \texttt{.gro}.
\end{itemize}

Por tanto, el PDB sirve únicamente para la \textbf{inspección visual inicial},
pero las distancias, ángulos y diedros se deben seleccionar siempre utilizando
los índices del archivo \texttt{.gro} o \texttt{.tpr}, que se muestran en la siguiente sección.

\subsubsection*{Inspección de los archivos \texttt{.gro} y \texttt{.top} con \texttt{less}}

Tras la ejecución de \texttt{gmx pdb2gmx}, se inspeccionaron los archivos generados 
(\texttt{arn.gro} y \texttt{arn.top}) utilizando el comando:

\begin{verbatim}
less arn.gro
less arn.top
\end{verbatim}

\paragraph{Estructura del archivo \texttt{.gro}.}
Cada línea del archivo \texttt{.gro} sigue el formato:

\begin{center}
\texttt{Residuo  NombreResiduo  NombreÁtomo  Índice  x  y  z}
\end{center}

Donde:
\begin{itemize}
  \item \textbf{Residuo}: número del residuo en la cadena (1 = ACE, 2 = ALA, 3 = ARG, 4 = ASN, 5 = NME).
  \item \textbf{NombreResiduo}: etiqueta de tres letras (ACE, ALA, ARG, ASN, NME).
  \item \textbf{NombreÁtomo}: nombre del átomo según CHARMM27 (N, HN, CA, HA, C, O…).
  \item \textbf{Índice}: número de átomo que GROMACS utiliza en todos los análisis.
  \item \textbf{x, y, z}: coordenadas cartesianas en nm.
\end{itemize}

Un fragmento representativo del archivo \texttt{arn\_box\_solv.gro} es:

\begin{verbatim}
    2ALA      N     7   1.044   1.361   1.264
    2ALA     HN     8   0.996   1.443   1.298
    2ALA     CA     9   1.185   1.368   1.225
    3ARG      N    17   1.378   1.412   1.364
    3ARG     HN    18   1.414   1.323   1.333
    3ARG     CA    19   1.453   1.494   1.458
\end{verbatim}

Este análisis permite verificar que:
\begin{itemize}
  \item 7--8--9 son N--HN--CA del residuo Ala (átomos consecutivos).
  \item 17--18--19 son N--HN--CA del residuo Arg.
  \item 9--17 representan un contacto corto CA(Ala)···N(Arg) entre residuos adyacentes.
  \item 17--18 forman el enlace covalente N--HN de Arg.
\end{itemize}

\paragraph{Estructura del archivo \texttt{.top}.}
El archivo \texttt{arn.top} contiene la descripción completa del sistema:

\begin{itemize}
  \item \textbf{Tipos atómicos} (columna \emph{type})
  \item \textbf{Número de residuo} (\emph{resnr})
  \item \textbf{Nombre del residuo} (\emph{residue})
  \item \textbf{Nombre del átomo} (\emph{atom})
  \item \textbf{Carga parcial} (\emph{charge})
  \item \textbf{Masa} (\emph{mass})
\end{itemize}

Un fragmento típico del bloque \texttt{[ atoms ]} es:

\begin{verbatim}
[ atoms ]
; nr  type  resnr residue  atom  cgnr charge   mass
  1   CT3      1   ACE     CH3    1   -0.27   12.011
  2    HA      1   ACE    HH31    2    0.09    1.008
  3    HA      1   ACE    HH32    3    0.09    1.008
  4    HA      1   ACE    HH33    4    0.09    1.008
  5     C      1   ACE       C    5    0.51   12.011
  6     O      1   ACE       O    6   -0.51   15.999
\end{verbatim}

Este bloque confirma:
\begin{itemize}
  \item que los grupos terminales ACE y NME son reconocidos correctamente,
  \item que las cargas parciales y tipos atómicos son los de CHARMM27,
  \item que la topología generada está preparada para simulación.
\end{itemize}

\paragraph*{Conclusión.}  
La inspección con \texttt{less} permite verificar que la estructura generada es 
correcta, que los índices atómicos usados en distancias/ángulos son válidos, y 
que la topología coincide con el campo de fuerza seleccionado.

\subsubsection*{Definición de la caja (\texttt{gmx editconf})}
\begin{lstlisting}[style=bashstyle]
# Caja cúbica de 3.0 nm de lado, soluto centrado
# -bt cubic: define la forma de la caja
# -box 3.0 3.0 3.0: dimensiones en nm (soluto centrado automáticamente)
gmx editconf -f arn.gro -o arn-box.gro -bt cubic -box 3.0 3.0 3.0
\end{lstlisting}
\textbf{Qué hace}. Define la celda periódica (no altera la topología) y fija el volumen para la solvatación.

\subsubsection*{Solvatación (\texttt{gmx solvate})}
\begin{lstlisting}[style=bashstyle]
# Llenar la caja con agua TIP3P y actualizar la topología
gmx solvate -cp arn-box.gro -cs -o arn_box_solv.gro -p arn.top

# (Opcional) PDB para inspección visual en PyMOL
gmx editconf -f arn_box_solv.gro -o arn_box_solv.pdb
\end{lstlisting}
\textbf{Qué hace}. Inserta moléculas de agua coherentes con TIP3P y actualiza \texttt{arn.top} con el número de SOL. El recuento típico total ronda \(\sim\)2640 átomos para este tamaño de caja.

\begin{lstlisting}[style=bashstyle]
pymol arn_box_solv.pdb
\end{lstlisting}

\begin{figure}[H]
  \centering
  \includegraphics[width=0.90\textwidth]{img/solvatacion.png}
  \caption{\textbf{Solvatación}. Péptido ARN centrado en caja cúbica de \SI{3.0}{nm} y rodeado de agua TIP3P.}
\end{figure}

\subsubsection*{Inspección del sistema solvatado}

Tras la solvatación, se verificó el contenido del archivo \texttt{arn\_box\_solv.gro}
mediante el comando:

\begin{verbatim}
less arn_box_solv.gro
\end{verbatim}

Un fragmento representativo del fichero es:

\begin{verbatim}
GRoups of Organic Molecules in ACtion for Science
 2640
    1ACE    CH3    1   0.837   1.239   1.294
    1ACE   HH31    2   0.825   1.167   1.374
    2ALA      N    7   1.044   1.361   1.264
    2ALA     HN    8   0.996   1.443   1.298
    2ALA     CA    9   1.185   1.368   1.225
...
  864SOL    HW1 2636   2.320   2.214   2.320
  864SOL    HW2 2637   2.251   2.246   2.465
  865SOL     OW 2638   1.901   2.939   2.162
  865SOL    HW1 2639   2.000   2.928   2.153
  865SOL    HW2 2640   1.861   2.853   2.194
   3.00000   3.00000   3.00000
\end{verbatim}

En este formato, cada línea corresponde a:

\begin{center}
\texttt{Residue  AtomName  AtomIndex  x  y  z}
\end{center}

Donde las coordenadas están en nanómetros.  
\paragraph*{Importante.} La presencia de los residuos \texttt{SOL} confirma la inserción de agua TIP3P.  
El encabezado indica que el sistema contiene \textbf{2640 átomos}, coherente con un
péptido pequeño rodeado de un volumen cúbico de agua de \(3.0\,\mathrm{nm}\).

Este análisis visual permite confirmar que:
\begin{itemize}
  \item el péptido aparece centrado en la caja,
  \item el agua llena completamente el volumen,
  \item no hay solapamientos extremos tras la construcción del sistema.
\end{itemize}

% ======================================================
% 3. BLOQUE DE EQUILIBRIO (EQUILIBRACIÓN NVT)
% ======================================================
\section{Bloque de EQUILIBRIO (Ensamble NVT)}
\label{sec:equilibrio}
\textbf{Objetivo}. Relajar el solvente y estabilizar la temperatura a \SI{298}{K} y \SI{400}{K} antes de la producción.

\subsection{Estructura de carpetas y ficheros (proyecto)}
\begin{lstlisting}[style=bashstyle]
# Carpeta raíz del proyecto
ARN_project=~/ARN_project

# Carpetas por temperatura (cada una con equilibration, run y analysis)
mkdir -p $ARN_project/298K/{equilibration,run,analysis}
mkdir -p $ARN_project/400K/{equilibration,run,analysis}
\end{lstlisting}

\paragraph{Copiar el archivo de control y el script SLURM.}
\begin{lstlisting}[style=bashstyle]
# 298 K
echo "Copiando archivos de control a 298K/equilibration"
cd $ARN_project/298K/equilibration
cp ~/for-students/equilibration-mdp/equiNVT.mdp .
cp ~/for-students/equilibration-mdp/submit_eck.sh .

# 400 K
echo "Copiando archivos de control a 400K/equilibration"
cd $ARN_project/400K/equilibration
cp ~/for-students/equilibration-mdp/equiNVT.mdp .
cp ~/for-students/equilibration-mdp/submit_eck.sh .
\end{lstlisting}

\paragraph{Copiar ficheros del sistema preparado.}
\begin{lstlisting}[style=bashstyle]
# 298 K
cd $ARN_project/298K/equilibration
cp $ARN_project/arn.top .
cp $ARN_project/arn_box_solv.gro .
cp $ARN_project/posre.itp .

# 400 K
cd $ARN_project/400K/equilibration
cp $ARN_project/arn.top .
cp $ARN_project/arn_box_solv.gro .
cp $ARN_project/posre.itp .
\end{lstlisting}

\subsection{Neutralización iónica (dentro del Equilibrio)}
\textbf{Motivación}. Aunque PME tolera carga neta con un fondo uniforme, neutralizar evita artefactos y representa condiciones realistas.

\subsubsection*{Paso 1: preparar \texttt{ion.tpr}}

\noindent\textbf{Nota}. gmx grompp = "preprocesador" de GROMACS. Combina parámetros (.mdp), coordenadas (.gro), topología (.top) y (si existe) checkpoint (.cpt) para generar un archivo .tpr que contiene TODA la configuración necesaria para ejecutar MD.

\begin{lstlisting}[style=bashstyle]
# .tpr auxiliar con parámetros de corte/PME para genion
gmx grompp -f equiNVT.mdp -c arn_box_solv.gro -p arn.top -o ion.tpr
\end{lstlisting}

\subsubsection*{Paso 2: insertar ion para neutralizar}

\noindent\textbf{Nota}. En lugar de añadir un número fijo de iones (como haría -nn), usamos la opción -neutral para que GROMACS calcule automáticamente cuántos iones necesita el sistema para que su carga neta sea exactamente cero.

\begin{lstlisting}[style=bashstyle]
# Sustituye una molécula de SOL por el ion necesario (p.ej., CL-)
gmx genion -s ion.tpr -p arn.top -o arn_neutral.gro -neutral
# Cuando pregunte por el grupo de solvente, seleccionar: SOL (número 13)
\end{lstlisting}

\paragraph*{Nota sobre el número total de átomos.}
Tras la solvatación, el sistema contiene \textbf{2640 átomos} (860 moléculas de
agua TIP3P más el péptido). Durante la neutralización iónica 
(\texttt{gmx genion -neutral}), GROMACS sustituye una molécula de agua TIP3P 
(compuesta por 3 átomos) por un ion cloruro Cl$^{-}$ (1 átomo). 
Como consecuencia, el total se reduce a \textbf{2638 átomos}, que es el valor 
empleado en todas las etapas posteriores de equilibración, producción y análisis.

\subsubsection*{Paso 3: verificaciones}
\begin{lstlisting}[style=bashstyle]
# Recompilar para comprobar ausencia de carga neta
gmx grompp -f equiNVT.mdp -c arn_neutral.gro -p arn.top -o arn_neutral.tpr
# Exportar PDB para inspección visual
gmx editconf -f arn_neutral.gro -o arn-neutral.pdb
# (Opcional) Distancia mínima ion-soluto
gmx mindist -f arn_neutral.gro -s arn_neutral.tpr -od mindist_ion.xvg
\end{lstlisting}

\subsection{Archivo de control de equilibración (\texttt{equiNVT.mdp})}
\begin{Verbatim}[fontsize=\small,frame=single]
; ===== Integradora y paso temporal =====
integrator = md           ; Integrador leap-frog
dt         = 0.0005       ; 0.5 fs

; ===== Duración del equilibrado =====
nsteps     = 800000       ; 400 ps (0.5 fs x 800k)

; ===== Tratamiento del centro de masas =====
constraints= none
comm-mode  = Linear
nstcomm    = 1000

; ===== Control de temperatura (NVT real) =====
Tcoupl     = v-rescale    ; Bussi thermostat 
tc-grps    = System       ; Acopla todo el sistema
tau_t      = 0.1          ; Respuesta rápida estable
ref_t      = 298          ; Usar 400 para la segunda eq

; ===== Velocidades iniciales =====
gen_vel    = yes
gen_temp   = 298          ; Usar 400 para la segunda eq
gen_seed   = 12345678

; ===== (Opcional) Restringir el soluto =====
; define     = -DPOSRES   ; Requiere posre.itp incluido en la .top
\end{Verbatim}
\textbf{Notas}. \texttt{v-rescale} garantiza muestreo canónico; \texttt{gen\_vel=yes} inicializa velocidades Maxwell–Boltzmann; \texttt{ref\_t}/\texttt{gen\_temp} son los únicos parámetros térmicos que cambian entre 298 y 400 K. También se fija explícitamente \texttt{nsteps = 800000}, que produce \textbf{400 ps}, la \emph{duración de equilibración} requerida por la práctica.

\subsection{Equilibración a 298 K}
\begin{lstlisting}[style=bashstyle]
cd $ARN_project/298K/equilibration
# Compilar y ejecutar
# Ajustar en equiNVT.mdp: ref_t = 298 y gen_temp = 298
gmx grompp -f equiNVT.mdp -c arn_neutral.gro -p arn.top -o arn_eq_298.tpr
bash submit_eck.sh # (mdrun -deffnm arn_eq_298 -nt 1)
gmx editconf -f arn_eq_298.gro -o arn_eq_298.pdb
\end{lstlisting}
\textbf{Salida esperada}. 400 ps completados; T media \(\approx\) \SI{298}{K} con fluctuaciones \(\pm\)8–9 K; redistribución del solvente sin colisiones, integración estable.

\begin{lstlisting}[style=bashstyle]
pymol arn_eq_298.pdb
\end{lstlisting}

\begin{figure}[H]
  \centering
  \includegraphics[width=0.90\textwidth]{img/equi_298.png}
  \caption{Equilibración a \SI{298}{K}.}
\end{figure}

\subsubsection*{Resumen cuantitativo de la equilibración (NVT)}

\begin{lstlisting}[style=bashstyle]
# Comprobar temperatura
# gmx energy lee el fichero .edr y permite extraer magnitudes termodinámicas registradas durante la simulación
# E. cinética, E. potencial, temperatura, presión...
# En el menú se deben seleccionar las variables por número o nombre.
gmx energy -f arn_eq_298.edr -o temperatura_298.xvg
\end{lstlisting}

\begin{table}[H]
\centering
\begin{tabular}{lcccc}
\toprule
\textbf{Modelo} & \textbf{T media (K)} & \textbf{Err.Est. (K)} & \textbf{RMSD (K)} & \textbf{Tot-Drift (K)} \\
\midrule
298 K (eq) & 298.104 & 0.25 & 8.866 & -1.589 \\
\bottomrule
\end{tabular}
\caption{Temperatura media durante la equilibración NVT (400 ps).}
\end{table}

\subsubsection*{Comparación estructural antes y después de la equilibración}

Para confirmar la estabilidad geométrica del sistema durante la fase NVT,
se compararon las primeras líneas del sistema neutralizado y del sistema
ya equilibrado:

\begin{verbatim}
head -n 10 arn_neutral.gro

Protein in water
 2638
    1ACE    CH3    1   0.837   1.239   1.294
    1ACE   HH31    2   0.825   1.167   1.374
    2ALA      N    7   1.044   1.361   1.264
    2ALA     HN    8   0.996   1.443   1.298
\end{verbatim}

Tras 400 ps de equilibración:

\begin{verbatim}
head -n 10 arn_eq_298.gro

Protein in water
 2638
    1ACE    CH3    1   1.344   1.494   1.258   0.5542 -0.3714  0.1199
    1ACE   HH31    2   1.370   1.497   1.371  -0.4920  1.3102  1.5642
    2ALA      N    7   1.503   1.651   1.192   0.4244 -0.5396  0.1994
    2ALA     HN    8   1.456   1.715   1.256  -1.6591  0.9063 -0.9113
\end{verbatim}

El formato de las líneas después de la equilibración incluye:

\begin{center}
\texttt{Residue AtomName Index  x y z   vx vy vz}
\end{center}

Que ahora incorpora las \textbf{velocidades atómicas} \((v_x, v_y, v_z)\)
generadas por el termostato.

Se observa que:
\begin{itemize}
  \item el número total de átomos permanece constante (2638),
  \item las posiciones cambian ligeramente debido al reordenamiento del solvente,
  \item las velocidades iniciales se han estabilizado según el reparto Maxwell–Boltzmann.
\end{itemize}

\subsection{Equilibración a 400 K}
\begin{lstlisting}[style=bashstyle]
cd $ARN_project/400K/equilibration
# Ajustar en equiNVT.mdp: ref_t = 400 y gen_temp = 400
gmx grompp -f equiNVT.mdp -c arn_neutral.gro -p arn.top -o arn_eq_400.tpr
bash submit_eck.sh   # (mdrun -deffnm arn_eq_400 -nt 1)
gmx editconf -f arn_eq_400.gro -o arn_eq_400.pdb
\end{lstlisting}
\textbf{Salida esperada}. 400 ps sin errores; T \(\approx\) \SI{400}{K} (RMSD \(\sim\)9 K); mayor amplitud de fluctuaciones, coherente con \(T\).

\begin{lstlisting}[style=bashstyle]
pymol arn_eq_400.pdb
\end{lstlisting}

\begin{figure}[H]
  \centering
  \includegraphics[width=0.90\textwidth]{img/equi_400.png}
  \caption{Equilibración a \SI{400}{K}.}
\end{figure}

\begin{lstlisting}[style=bashstyle]
# Comprobar temperatura
# gmx energy lee el fichero .edr y permite extraer magnitudes termodinámicas registradas durante la simulación
# E. cinética, E. potencial, temperatura, presión...
# En el menú se deben seleccionar las variables por número o nombre.
gmx energy -f arn_eq_400.edr -o temperatura_400.xvg
\end{lstlisting}

\subsubsection*{Resumen cuantitativo de la equilibración (NVT)}

\begin{table}[H]
\centering
\begin{tabular}{lcccc}
\toprule
\textbf{Modelo} & \textbf{T media (K)} & \textbf{Err.Est. (K)} & \textbf{RMSD (K)} & \textbf{Tot-Drift (K)} \\
\midrule
400 K (eq) & 400.185 & 0.25 & 9.157 & -1.267 \\
\bottomrule
\end{tabular}
\caption{Temperatura media durante la equilibración NVT (400 ps).}
\end{table}

\subsubsection*{Comparación estructural a 298 K y 400 K}

De forma análoga, se verificó la estructura inicial del sistema a 400 K:

\begin{verbatim}
head -n 10 arn_eq_400.gro

Protein in water
 2638
    1ACE    CH3    1   0.387   2.560   1.372   0.0911  0.4564  0.3349
    1ACE   HH31    2   0.417   2.511   1.282  -2.0350  0.9370 -1.3482
    2ALA      N    7   0.521   2.725   1.505  -0.1331  0.4494  1.0357
    2ALA     HN    8   0.507   2.672   1.585  -0.8350 -0.8074 -1.5644
\end{verbatim}

Los valores de velocidad son mayores que en el caso de 298 K, lo que refleja
la relación esperada entre energía cinética y temperatura:

\[
E_k \propto T.
\]

Además, los átomos de agua muestran desplazamientos algo mayores, coherentes con
la mayor movilidad térmica del sistema.

% ======================================================
% 4. BLOQUE de RUN (PRODUCCIÓN NVT)
% ======================================================
\section{Bloque de RUN (Producción NVT)}
\label{sec:run}
\textbf{Objetivo}. Generar trayectorias de producción con \textbf{continuidad física} desde el equilibrado (no se regeneran velocidades).

\subsection{Preparación del entorno de ejecución}
\begin{lstlisting}[style=bashstyle]
# 298 K
cd $ARN_project/298K/run
cp ../equilibration/arn.top .
cp ../equilibration/arn_eq_298.gro .
cp ../equilibration/arn_eq_298.cpt .
cp ../equilibration/posre.itp .

# 400 K
cd $ARN_project/400K/run
cp ../equilibration/arn.top .
cp ../equilibration/arn_eq_400.gro .
cp ../equilibration/arn_eq_400.cpt .
cp ../equilibration/posre.itp .
\end{lstlisting}
\textbf{Resumen técnico}. \texttt{arn\_eq\_XXX.gro} (estructura final), \texttt{arn\_eq\_XXX.cpt} (checkpoint con velocidades), \texttt{posre.itp} (si se activa POSRES), \texttt{arn.top} (topología coherente).

\subsection{Archivo de control de RUN (\texttt{runNVT.mdp})}
\begin{Verbatim}[fontsize=\small,frame=single]
; ===== Integración =====
integrator   = md

; ===== Paso y duración =====
dt           = 0.0005    ; 0.5 fs
nsteps       = 4000      ; 2 ps

; ===== Homogeneidad con la eq =====
constraints  = none
comm-mode    = Linear
nstcomm      = 1000

; ===== Termostato =====
Tcoupl       = v-rescale
tc-grps      = System
tau_t        = 0.1
ref_t        = 298       ; usar 400 para la segunda producción

; ===== Velocidades =====
gen_vel      = no        ; continuidad (se leen del .cpt)
\end{Verbatim}
\textbf{Notas}. \texttt{gen\_vel=no} garantiza continuidad dinámica; \texttt{nsteps = 4000} produce \textbf{2 ps}, la \emph{duración de producción} requerida.

\subsection{Generación del \texttt{.tpr}}
\begin{lstlisting}[style=bashstyle]
# 298 K
gmx grompp -f runNVT.mdp -c arn_eq_298.gro -t arn_eq_298.cpt -p arn.top -o arn_run_298.tpr
# 400 K
gmx grompp -f runNVT.mdp -c arn_eq_400.gro -t arn_eq_400.cpt -p arn.top -o arn_run_400.tpr
\end{lstlisting}
\textbf{Validación}. El preprocesador confirma lectura de coordenadas/velocidades, grados de libertad y malla PME sin inconsistencias topológicas.

\subsection{Ejecución con SLURM}
\noindent\textbf{Nota}. gmx mdrun = motor de dinámica molecular. Lee el .tpr, calcula fuerzas, integra ecuaciones de Newton y produce .trr, .edr, .log, .gro final.
\begin{lstlisting}[style=bashstyle]
#!/bin/bash
#SBATCH -p eck-q
#SBATCH --chdir=$ARN_project/298K/run
#SBATCH -J run298
#SBATCH --cpus-per-task=1

date
gmx mdrun -deffnm arn_run_298 -nt 1
date

# Ejecución equivalente para 400 K:
# gmx mdrun -deffnm arn_run_400 -nt 1
\end{lstlisting}
\textbf{Salida esperada}. Lectura correcta del \texttt{.tpr}, ejecución sin errores, \(\sim\)2000 \emph{frames} (2 ps), rendimiento \(\sim\)21 ns/día y ficheros: \texttt{.trr}, \texttt{.edr}, \texttt{.log}, \texttt{.tpr}.

\subsection{Exportación para visualización}
\begin{lstlisting}[style=bashstyle]
# Sistema completo (Protein + agua + ion -> número 0)
gmx traj -f arn_run_298.trr -s arn_run_298.tpr -oxt cartoon.pdb
# Solo soluto (seleccionar 'Protein' -> número 1)
gmx traj -f arn_run_298.trr -s arn_run_298.tpr -oxt arn_solo.pdb
\end{lstlisting}

\noindent\textbf{Nota.} En la visualización del \emph{sistema completo} en PyMOL , la opción \texttt{Display -> Roving Detail} permite resaltar automáticamente detalles locales del entorno (geometría, contactos y posibles enlaces de hidrógeno), facilitando la inspección de la estructura en movimiento.

\begin{lstlisting}[style=bashstyle]
# Sistema completo (Protein + agua + ion)
pymol cartoon.pdb
# Solo soluto
pymol arn_solo.pdb
\end{lstlisting}

% --- Figura: Sistema completo 298 K ---
\begin{figure}[H]
  \centering
  \includegraphics[width=0.90\textwidth]{img/arn_298_completo.png}
  \caption{\textbf{298 K – Sistema completo}. Representación del tripéptido ARN junto con el solvente e ion, visualizado con la opción \textit{Roving Detail}.}
\end{figure}

% --- Figura: Solo soluto 298 K ---
\begin{figure}[H]
  \centering
  \includegraphics[width=0.90\textwidth]{img/arn_solo_298.png}
  \caption{\textbf{298 K – Soluto únicamente}. Visualización del tripéptido ARN sin agua ni iones, lo que permite observar una conformación compacta y estable a esa temperatura.}
\end{figure}

\begin{lstlisting}[style=bashstyle]
# Sistema completo (Protein + agua + ion)
gmx traj -f arn_run_400.trr -s arn_run_400.tpr -oxt cartoon.pdb
# Solo soluto (seleccionar 'Protein')
gmx traj -f arn_run_400.trr -s arn_run_400.tpr -oxt arn_solo.pdb
\end{lstlisting}

% --- Figura: Sistema completo 400 K ---
\begin{figure}[H]
  \centering
  \includegraphics[width=0.90\textwidth]{img/arn_400_completo.png}
  \caption{\textbf{400 K – Sistema completo}. El incremento térmico provoca un movimiento más amplio y fluctuaciones intensas del tripéptido ARN en el seno del disolvente (visualizado con \textit{Roving Detail}).}
\end{figure}

% --- Figura: Solo soluto 400 K ---
\begin{figure}[H]
  \centering
  \includegraphics[width=0.90\textwidth]{img/run_400_solo.png}
  \caption{\textbf{400 K – Soluto únicamente}. Representación del tripéptido ARN sin agua ni iones. Una mayor energía térmica genera una conformación más extendida, con mayor movilidad torsional y desplazamientos amplios del esqueleto peptídico.}
\end{figure}

% ======================================================
% 5. BLOQUE DE ANÁLISIS (POSTPROCESADO)
% ======================================================
\section{Bloque de ANÁLISIS (Postprocesado)}
\label{sec:analisis}

\noindent
\textbf{Objetivo}. Cuantificar propiedades estructurales y energéticas a partir de los archivos \texttt{.trr}, \texttt{.edr} y \texttt{.tpr}.

\noindent
\textbf{Escalas y unidades}. Distancias (nm), ángulos (°), energía (kJ/mol), temperatura (K), tiempo (ps), velocidades (nm/ps).

\noindent
\textbf{Figuras}. Todas las figuras incluidas en este informe fueron generadas con \textbf{Python} a partir de los archivos \texttt{.xvg} producidos por GROMACS, utilizando las bibliotecas \texttt{numpy}, \texttt{pandas} y \texttt{matplotlib}.
% ======================================================
\subsection{Preparación del directorio de análisis}

\begin{lstlisting}[style=bashstyle]
# 298 K
# 298 K — copio todos los archivos generados durante la producción NVT
# .tpr -> topología + parámetros
# .gro -> última estructura
# .trr -> trayectoria completa
# .edr -> energías
# .log -> registro de ejecución
cp ../run/arn_run_298.tpr .
cp ../run/arn_run_298.gro .
cp ../run/arn_run_298.trr .
cp ../run/arn_run_298.edr .
cp ../run/arn_run_298.log .

# 400 K — lo mismo para la producción a temperatura alta
cp ../run/arn_run_400.tpr .
cp ../run/arn_run_400.gro .
cp ../run/arn_run_400.trr .
cp ../run/arn_run_400.edr .
cp ../run/arn_run_400.log .
\end{lstlisting}

Los archivos obtenidos permiten realizar todos los análisis.

% ======================================================
\subsection{Radio de giro}

\noindent\textbf{Definición.} El radio de giro (\(R_{\mathrm{g}}\)) se define como:

\[
R_{\mathrm{g}}^{2} = \frac{1}{M} \sum_{i=1}^{N} m_i \left\lVert \mathbf{r}_i - \mathbf{R}_{\mathrm{CM}} \right\rVert^2,
\]

Donde \(M\) es la masa total, \(m_i\) la masa del átomo \(i\), \(\mathbf{r}_i\) su posición y 
\(\mathbf{R}_{\mathrm{CM}}\) el centro de masas.

\noindent\textbf{Nota}. En la ejecución interactiva de gmx gyrate se seleccionó el grupo 'Protein' -> número 1, que en este sistema corresponde normalmente a la opción 1 del menú de selección.

\begin{lstlisting}[style=bashstyle]
# Calcula el radio de giro (Rg) a partir de la trayectoria
# -f: trayectoria .trr   -s: topología   -o: salida .xvg
gmx gyrate -f arn_run_298.trr -s arn_run_298.tpr -o gyrate_298.xvg
gmx gyrate -f arn_run_400.trr -s arn_run_400.tpr -o gyrate_400.xvg
\end{lstlisting}

\begin{figure}[H]
\centering
\includegraphics[width=0.90\textwidth]{img/figura_python_gyrate_298K.png}
\caption{Radio de giro a 298 K.}
\end{figure}

\begin{figure}[H]
\centering
\includegraphics[width=0.90\textwidth]{img/figura_python_gyrate_400K.png}
\caption{Radio de giro a 400 K.}
\end{figure}

\paragraph*{Análisis del radio de giro}

Las Figuras anteriores muestran la evolución temporal del radio de giro (Rg), una
magnitud que cuantifica la compacidad global del tripéptido. En ambos regímenes
térmicos, el Rg presenta oscilaciones suaves y continuas, reflejo del
desplazamiento colectivo de los átomos del backbone y de la reorganización del
entorno inmediato.

A \SI{298}{K}, el radio de giro se mantiene en torno a
\(0.435{-}0.453~\mathrm{nm}\), indicando una conformación relativamente compacta y
estable. A \SI{400}{K}, el Rg disminuye ligeramente hasta valores próximos a
\(0.398{-}0.424~\mathrm{nm}\), mostrando mayores fluctuaciones y episodios de mayor
contracción estructural. Este comportamiento es coherente con un incremento en la
flexibilidad interna inducido por la temperatura, que favorece la exploración de
estados conformacionales más dinámicos y menos extendidos.

Sin embargo, la duración de la simulación es demasiado breve para extraer
conclusiones estadísticas sólidas: los datos no están decorrelacionados y no
puede evaluarse un promedio convergente ni intervalos de confianza. Por tanto,
la aparente disminución del Rg a \SI{400}{K} debe interpretarse únicamente como
una \emph{tendencia observada} en esta ventana temporal, no como un cambio
estructural concluyente.

\subsection{Verificación de índices atómicos}

Antes de seleccionar los átomos para distancias y ángulos, se inspeccionó
el archivo \texttt{arn\_box\_solv.gro} utilizando:

\begin{verbatim}
less arn_box_solv.gro
\end{verbatim}

Las líneas relevantes del fichero son:

\begin{verbatim}
    2ALA      N      7
    2ALA     HN      8
    2ALA     CA      9
    2ALA      C     15
    2ALA      O     16
    ...
    3ARG      N     17
    3ARG     HN     18
    3ARG     CA     19
\end{verbatim}

Estos índices demuestran que:
\begin{itemize}
  \item \textbf{7--9--15} son átomos consecutivos del backbone de ALA-2 (N--CA--C),
  \item \textbf{9--15--17} conectan el C de ALA-2 con el N de ARG-3 formando el \textit{enlace peptídico},
  \item \textbf{7--9--15--17} definen el diedro N--CA--C--N (\(\psi\) del residuo ALA-2),
  \item \textbf{9--15--17--19} definen el diedro CA--C--N--CA (\(\omega\) del enlace peptídico entre ALA-2 y ARG-3).
\end{itemize}

Con esta verificación se garantiza que todos los átomos seleccionados 
para distancias, ángulos y diedros cumplen el requisito de la práctica:
\emph{“los átomos deben estar juntos o ser contiguos en el esqueleto peptídico”}.

\noindent\textbf{Nota previa}. Antes de definir el archivo \texttt{distancias.ndx}, se verificaron visualmente en PyMOL los átomos implicados en cada enlace. Para ello se emplearon las siguientes selecciones:

\begin{verbatim}
pymol select enlace1, id 7+9     # N–CA (ALA-2)
pymol select enlace2, id 9+15    # CA–C (ALA-2)
\end{verbatim}

\begin{figure}[H]
    \centering

    % --- Imagen 1 ---
    \includegraphics[width=0.60\textwidth]{pymol_enlace1.png}
    \caption*{\small Selección en PyMOL del enlace 7--9 (\texttt{select enlace1, id 7+9 }).}

    \vspace{0.4cm}

    % --- Imagen 2 ---
    \includegraphics[width=0.70\textwidth]{pymol_enlace2.png}
    \caption*{\small Selección en PyMOL del enlace 9--15 (\texttt{select enlace2, id 9+15}).}

    \caption{\textbf{Verificación visual en PyMOL} de los pares atómicos usados para el análisis de distancias.}
\end{figure}

Esto permitió confirmar los pares atómicos seleccionados para el análisis 
geométrico del backbone. En concreto, se identificaron los enlaces consecutivos del 
residuo ALA-2: el par \textbf{7--9} corresponde al enlace covalente \textbf{N--CA}, mientras que el par 
\textbf{9--15} define el enlace covalente \textbf{CA--C}. Estos dos enlaces, adyacentes en el esqueleto 
peptídico, constituyen las unidades estructurales básicas del backbone y representan magnitudes idóneas 
para evaluar su estabilidad geométrica bajo distintas condiciones de temperatura.

\noindent\textbf{Identificación de los enlaces analizados.} 
El par de átomos 7--9 corresponde al \textbf{enlace covalente N--CA} del residuo ALA-2. 
El par 9--15 corresponde al \textbf{enlace covalente CA--C} del mismo residuo. 
Ambos pares representan enlaces consecutivos del backbone y son ideales para evaluar 
la estabilidad geométrica del esqueleto peptídico bajo diferentes condiciones de temperatura.

\subsection{Distancias interatómicas}

\noindent\textbf{Nota}. Al ejecutar gmx distance, GROMACS solicita seleccionar el grupo definido en el archivo distancias.ndx. Se seleccionó \textbf{Opción 0: \texttt{[ enlace\_N\_CA ]}} seguida de Enter para luego selecciona \textbf{Opción 1: \texttt{[ enlace\_CA\_C ]}} para calcular simultáneamente ambos enlaces del backbone.

\begin{lstlisting}[style=bashstyle]
# Creo un archivo .ndx con los pares de átomos que quiero medir
# 7-9 = enlace N-CA ; 9-15 = enlace CA-C
cat > dist_backbone.ndx << 'EOF'
[ enlace_N_CA ]
7 9
[ enlace_CA_C ]
9 15
EOF

# Distancias a 298 K -> ficheros nuevos
gmx distance -f arn_run_298.trr -s arn_run_298.tpr -n dist_backbone.ndx -oall dist_backbone_298.xvg

# Distancias a 400 K -> ficheros nuevos
gmx distance -f arn_run_400.trr -s arn_run_400.tpr -n dist_backbone.ndx -oall dist_backbone_400.xvg
\end{lstlisting}

\subsubsection*{Análisis de las distancias internas}

Los siguientes valores fueron proporcionados directamente por \texttt{gmx distance}
a partir de las trayectorias simuladas. Se analizaron los enlaces consecutivos
\textbf{N--CA} y \textbf{CA--C}, que comparten el átomo \textbf{CA}.

\begin{table}[h!]
    \centering
    \caption{Distancias promedio y desviaciones estándar obtenidas con GROMACS
    para los enlaces del backbone a 298 K y 400 K.}
    \begin{tabular}{lccc}
        \hline
        \textbf{Enlace} & \textbf{Temperatura} & \textbf{Distancia media (nm)} & \textbf{Desv. estándar (nm)} \\
        \hline
        N--CA & 298 K & 0.14516 & 0.00228 \\
        N--CA & 400 K & 0.14539 & 0.00460 \\
        CA--C & 298 K & 0.15280 & 0.00288 \\
        CA--C & 400 K & 0.15288 & 0.00443 \\
        \hline
    \end{tabular}
\end{table}

Las distancias covalentes N--CA y CA--C muestran valores medios prácticamente
idénticos entre \SI{298}{K} y \SI{400}{K}, lo que indica que el aumento de
temperatura no altera su longitud promedio. No obstante, las fluctuaciones se
incrementan de forma clara a \SI{400}{K}, duplicándose aproximadamente las
desviaciones estándar. Este comportamiento es coherente con el aumento de la
energía térmica: los enlaces mantienen su geometría media, pero vibran con mayor
amplitud.

\begin{figure}[H]
\centering
\includegraphics[width=0.90\textwidth]{img/figura_python_distancias_298K.png}
\caption{Distancias de backbone N–CA (7--9) y CA–C (9--15) a 298 K.}
\end{figure}

\begin{figure}[H]
\centering
\includegraphics[width=0.90\textwidth]{img/figura_python_distancias_400K.png}
\caption{Distancias de backbone N–CA (7--9) y CA–C (9--15) a 400 K.}
\end{figure}

\paragraph*{Comparación de las distancias de enlace a 298 K y 400 K}

La comparación directa entre ambas figuras revela que las distancias del backbone
\textbf{N--CA (7--9)} y \textbf{CA--C (9--15)} se mantienen extraordinariamente estables 
tanto a \SI{298}{K} como a \SI{400}{K}, oscilando alrededor de sus valores medios sin 
mostrar derivas apreciables a lo largo de los \SI{2}{ps} de simulación. En la gráfica de 
\SI{298}{K}, ambas señales presentan fluctuaciones muy estrechas y regulares, con un 
perfil prácticamente uniforme durante todo el intervalo temporal. En la gráfica de 
\SI{400}{K} se observa un incremento moderado en la amplitud de las oscilaciones, 
reflejo de la mayor agitación térmica, pero las trayectorias permanecen igualmente 
acotadas y centradas en los mismos valores medios que a \SI{298}{K}. En conjunto, la 
comparación demuestra que el aumento de temperatura no altera las longitudes 
medias de los enlaces covalentes del esqueleto peptídico, limitándose únicamente a 
amplificar las fluctuaciones instantáneas sin comprometer la integridad estructural 
del backbone.

%======================================================
\subsection{Ángulos internos}

\noindent\textbf{Nota}. En el menú de \texttt{gmx angle} se seleccionaron los grupos definidos en
\texttt{ang\_backbone.ndx}:
Opción 0: \texttt{[ ang\_N\_CA\_C ]} (átomos 7--9--15, ángulo N--CA--C),\\
Opción 1: \texttt{[ ang\_CA\_C\_N ]} (átomos 9--15--17, ángulo CA--C--N).

\begin{lstlisting}[style=bashstyle]
# Creo archivo .ndx con los ángulos elegidos: 7-9-15 y 9-15-17
echo -e "[ ang_N_CA_C ]\n 7 9 15\n[ ang_CA_C_N ]\n 9 15 17" > ang_backbone.ndx

# 298 K (dos ejecuciones, un grupo cada vez)
# Para el primer ángulo (Selecciona el grupo 0 cuando lo pida)
gmx angle -f arn_run_298.trr -n ang_backbone.ndx -ov ang_N_CA_C_298.xvg -xvg none
# Para el segundo ángulo (Selecciona el grupo 1 cuando lo pida)
gmx angle -f arn_run_298.trr -n ang_backbone.ndx -ov ang_CA_C_N_298.xvg -xvg none

# 400 K
# Para el primer ángulo (Selecciona el grupo 0 cuando lo pida)
gmx angle -f arn_run_400.trr -n ang_backbone.ndx -ov ang_N_CA_C_400.xvg -xvg none
# Para el segundo ángulo (Selecciona el grupo 1 cuando lo pida)
gmx angle -f arn_run_400.trr -n ang_backbone.ndx -ov ang_CA_C_N_400.xvg -xvg none
\end{lstlisting}

\subsubsection*{Análisis de ángulos internos}

Los dos ángulos analizados corresponden al esqueleto peptídico:  
(i) \textbf{N--CA--C (7--9--15)},  
(ii) \textbf{CA--C--N (9--15--17)}.

\begin{table}[H]
\centering
\begin{tabular}{lcc}
\toprule
\textbf{Ángulo} & \textbf{298 K (°)} & \textbf{400 K (°)} \\
\midrule
N--CA--C (7--9--15) & 110.37 & 116.05 \\
CA--C--N (9--15--17) & 116.97 & 118.28 \\
\bottomrule
\end{tabular}
\caption{Valores medios de los ángulos internos del backbone a 298 K y 400 K.}
\end{table}

\noindent\textbf{Interpretación}.  
Los valores medios muestran que ambos ángulos mantienen una geometría estable en las dos temperaturas, con incrementos moderados a \SI{400}{K} que no alteran la estructura básica del backbone. El ángulo N--CA--C aumenta de \(\sim 110^\circ\) a \(\sim 116^\circ\), mientras que CA--C--N pasa de \(\sim 117^\circ\) a \(\sim 118^\circ\), variaciones compatibles con una mayor flexibilidad térmica. Las desviaciones estándar (\(2{-}4^\circ\)) indican fluctuaciones acotadas y coherentes con el comportamiento esperable de un esqueleto covalente rígido sometido a dinámica molecular en régimen NVT.

\begin{figure}[H]
\centering
\includegraphics[width=0.90\textwidth]{img/figura_python_angulos_298K.png}
\caption{Ángulos internos a 298 K.}
\end{figure}

\begin{figure}[H]
\centering
\includegraphics[width=0.90\textwidth]{img/figura_python_angulos_400K.png}
\caption{Ángulos internos a 400 K.}
\end{figure}

\paragraph*{Comparación de los ángulos de enlace}

Las figuras comparan la evolución temporal de los ángulos \textbf{N--CA--C (7--9--15)} y 
\textbf{CA--C--N (9--15--17)} del backbone a \SI{298}{K} y \SI{400}{K}. En ambos regímenes térmicos, los dos
ángulos oscilan alrededor de valores medios bien definidos, sin mostrar derivas apreciables durante los 
\SI{2}{ps} de simulación. A \SI{298}{K}, las fluctuaciones son más estrechas y regulares: el ángulo N--CA--C 
permanece predominantemente cercano a \(110^\circ\), mientras que el ángulo CA--C--N se mantiene alrededor de 
\(117{-}118^\circ\), describiendo variaciones moderadas y suavemente acotadas.

A \SI{400}{K}, ambos ángulos conservan prácticamente las mismas posiciones medias, pero las trayectorias 
presentan oscilaciones más amplias y frecuentes, reflejo directo del aumento de la movilidad térmica. Aun así, 
las señales permanecen confinadas dentro de sus intervalos característicos y no se observan cambios estructurales 
relevantes ni desviaciones persistentes respecto a sus valores centrales. En conjunto, la comparación muestra que 
el incremento de temperatura amplifica las fluctuaciones instantáneas sin alterar la geometría promedio del 
esqueleto peptídico, evidenciando la estabilidad de los ángulos de enlace N--CA--C y CA--C--N incluso bajo 
condiciones térmicas elevadas.
% ======================================================
\subsection{Ángulos diedros}

\noindent\textbf{Nota sobre los diedros analizados}. 
Los diedros seleccionados, \textbf{N--CA--C--N (7--9--15--17)} y 
\textbf{CA--C--N--CA (9--15--17--19)}, corresponden a torsiones internas consecutivas del 
backbone peptídico. Estos ángulos describen la rotación en torno a los enlaces 
C$\alpha$--C y C--N del esqueleto, pero \textbf{no} coinciden con los ángulos 
\(\phi\) y \(\psi\) clásicos del mapa de Ramachandran. Por tanto, sus valores no 
deben compararse directamente con \(\phi/\psi\).

\noindent\textbf{Nota}. En el menú interactivo de \texttt{gmx angle} (modo \textit{dihedral}) se seleccionaron:
opción 0: \texttt{[ dih\_N\_CA\_C\_N ]} (7--9--15--17),\quad
opción 1: \texttt{[ dih\_CA\_C\_N\_CA ]} (9--15--17--19).

\begin{lstlisting}[style=bashstyle]
# Archivo para diedros de backbone
echo -e "[ dih_N_CA_C_N ]\n 7 9 15 17\n[ dih_CA_C_N_CA ]\n 9 15 17 19" > dih_backbone.ndx

# 298 K
# Diedro 1: 7-9-15-17  (N–CA–C–N)
echo 0 | gmx angle -f arn_run_298.trr -n dih_backbone.ndx \
  -type dihedral -ov dih_N_CA_C_N_298.xvg -xvg none

# Diedro 2: 9-15-17-19 (CA–C–N–CA)
echo 1 | gmx angle -f arn_run_298.trr -n dih_backbone.ndx \
  -type dihedral -ov dih_CA_C_N_CA_298.xvg -xvg none

  # 400 K
# Diedro 1: 7-9-15-17  (N–CA–C–N)
echo 0 | gmx angle -f arn_run_400.trr -n dih_backbone.ndx \
  -type dihedral -ov dih_N_CA_C_N_400.xvg -xvg none

  # Diedro 2: 9-15-17-19 (CA–C–N–CA)
echo 1 | gmx angle -f arn_run_400.trr -n dih_backbone.ndx \
  -type dihedral -ov dih_CA_C_N_CA_400.xvg -xvg none
\end{lstlisting}

\subsubsection*{Análisis de ángulos diedros del esqueleto peptídico}

Los dos diedros analizados corresponden a torsiones internas consecutivas del backbone peptídico. Ninguno de ellos corresponde a los ángulos de Ramachandran
(\(\phi,\psi\)), por lo que no deben interpretarse como tales.:  
(i) \textbf{N--CA--C--N (7--9--15--17)},  
(ii) \textbf{CA--C--N--CA (9--15--17--19)}.

\noindent\textbf{Resultados de GROMACS}.  
El comando \texttt{gmx angle -type dihedral} proporciona, además de los valores
instantáneos, una \emph{media lineal} del ángulo y el parámetro de orden \(S^{2}\).
En variables periódicas como los diedros, la media lineal puede resultar
engañosa: valores próximos a \(+180^\circ\) y a \(-180^\circ\) se cancelan,
produciendo promedios artificiales cercanos a \(0^\circ\). El propio \(S^{2}\),
sin embargo, sí refleja la rigidez real del ángulo.

\begin{table}[H]
\centering
\small
\begin{tabular}{lcccccc}
\toprule
\multirow{2}{*}{\textbf{Diedro}} 
& \multicolumn{3}{c}{\textbf{298 K}} 
& \multicolumn{3}{c}{\textbf{400 K}} \\
\cmidrule(lr){2-4}\cmidrule(lr){5-7}
& \textbf{Media (°)} & \textbf{Std (°)} & $\mathbf{S^{2}}$ 
& \textbf{Media (°)} & \textbf{Std (°)} & $\mathbf{S^{2}}$ \\
\midrule
N--CA--C--N (7--9--15--17)    
&   9.72  &  8.73   & 0.977  
& $-42.33$ & 11.40  & 0.961 \\
CA--C--N--CA (9--15--17--19)  
& $-32.25$ & 169.98 & 0.979  
&   55.37  & 165.17 & 0.986 \\
\bottomrule
\end{tabular}
\caption{Medias lineales, desviaciones estándar y parámetros de orden $S^2$ de los dos diedros analizados. Las medias lineales no deben interpretarse como orientaciones reales debido a la periodicidad angular; el valor de $S^2$ es el indicador fiable de rigidez torsional.}
\end{table}

\noindent\textbf{Interpretación}.  
El diedro \textbf{N--CA--C--N} presenta una banda estrecha de valores en ambos
regímenes térmicos, lo que indica que la torsión alrededor del enlace C$\alpha$--C es
flexible pero acotada. A 400 K se observa un desplazamiento hacia valores
negativos (aprox. \(-40^\circ\)), consistente con una mayor exploración
conformacional inducida por la temperatura.

El diedro \textbf{CA--C--N--CA}, que corresponde al ángulo \(\omega\) del enlace
peptídico, aparece con medias lineales próximas a \(0^\circ\); sin embargo, esto es
un artefacto matemático. El parámetro de orden \(S^{2}\approx 0.97{-}0.99\) indica
una rigidez muy alta, compatible con un enlace peptídico en conformación
\textbf{trans} (\(\pm 180^\circ\)) durante toda la simulación. La desviación estándar
tan grande (\(\sim 165^\circ\)) simplemente refleja el cruce numérico entre \(-180^\circ\)
y \(+180^\circ\).

En resumen, el incremento térmico a 400 K aumenta la flexibilidad del diedro
\textbf{N--CA--C--N}, mientras que el ángulo \(\omega\) permanece rígido y
preferentemente trans, como es esperable en un enlace peptídico. Las medias
lineales deben interpretarse con cautela y no representan el valor físico real
del diedro debido a la periodicidad angular.

\begin{figure}[H]
\centering
\includegraphics[width=0.90\textwidth]{img/figura_python_diedros_tiempo_298K.png}
\caption{Ángulos diedros a 298 K.}
\end{figure}

\begin{figure}[H]
\centering
\includegraphics[width=0.90\textwidth]{img/figura_python_diedros_tiempo_400K.png}
\caption{Ángulos diedros a 400 K.}
\end{figure}

\paragraph*{Comparación de los ángulos diedros (298 K vs 400 K)}
En la figura de \SI{298}{K}, el diedro \textbf{N--CA--C--N (7--9--15--17)} evoluciona en una banda estrecha
próxima a \(0^\circ\), con oscilaciones moderadas y sin saltos, mientras que el diedro \textbf{CA--C--N--CA
(9--15--17--19)} permanece mayoritariamente en torno a \(\pm 180^\circ\), con cambios bruscos que reflejan
el cruce periódico entre los estados equivalentes trans.

A \SI{400}{K}, \textbf{N--CA--C--N} muestra un \emph{desplazamiento claro del valor medio} hacia
ángulos negativos (alrededor de \(-60^\circ\)) y un \emph{aumento de la amplitud} de las oscilaciones,
indicando mayor flexibilidad torsional. Por el contrario, \textbf{CA--C--N--CA} sigue confinado en el
entorno trans (\(\pm 180^\circ\)), aunque con \emph{transiciones más frecuentes} entre las dos ramas
debido a la periodicidad del ángulo. En conjunto, el incremento térmico amplifica la dinámica del diedro
\textbf{N--CA--C--N}, mientras que mantiene \textbf{CA--C--N--CA (}\(\omega\)\textbf{)}
rígido en trans, con cruces periódicos sin pérdida de estabilidad estructural.

\subsection{Mapas de Ramachandran}

\noindent
Para caracterizar la flexibilidad torsional del péptido, se calcularon los ángulos
\(\phi\) y \(\psi\) de cada residuo mediante el comando \texttt{gmx rama}. Este análisis
proporciona directamente un archivo \texttt{.xvg} con los valores de \(\phi/\psi\) por
frame y permite representar los mapas de Ramachandran individuales para
cada residuo del tripéptido. Aunque el sistema es pequeño y no presenta
estructura secundaria definida, el mapa \(\phi/\psi\) ofrece información clara sobre
las regiones conformacionales accesibles a distintas temperaturas y sobre la
sensibilidad torsional de cada residuo.

\begin{lstlisting}
gmx rama -s arn_run_298.tpr -f arn_run_298.trr -o rama_298.xvg -xvg none
# rama: -s .tpr (referencia) -f .trr (trayectoria) -> -o .xvg (phi/psi)
# -xvg none: salida limpia para procesar con grep/awk
# Salida con etiquetas del tipo: ALA-2  phi  psi; ARG-3  phi  psi; ASN-4  phi  psi
gmx rama -s arn_run_400.tpr -f arn_run_400.trr -o rama_400.xvg -xvg none
\end{lstlisting}

\begin{figure}[H]
\centering
\includegraphics[width=0.32\textwidth]{img/ramachandran_ALA-2_298K.png}
\includegraphics[width=0.32\textwidth]{img/ramachandran_ARG-3_298K.png}
\includegraphics[width=0.32\textwidth]{img/ramachandran_ASN-4_298K.png}
\caption{\textbf{Mapas de Ramachandran a \SI{298}{K}} para los residuos ALA-2, ARG-3 y ASN-4.}
\end{figure}

\begin{figure}[H]
\centering
\includegraphics[width=0.32\textwidth]{img/ramachandran_ALA-2_400K.png}
\includegraphics[width=0.32\textwidth]{img/ramachandran_ARG-3_400K.png}
\includegraphics[width=0.32\textwidth]{img/ramachandran_ASN-4_400K.png}
\caption{\textbf{Mapas de Ramachandran a \SI{400}{K}} para los residuos ALA-2, ARG-3 y ASN-4.}
\end{figure}

\paragraph*{Comparación de los mapas de Ramachandran a 298 K y 400 K}

La comparación entre los mapas de Ramachandran a \SI{298}{K} y \SI{400}{K} revela
cambios claramente dependientes del residuo. En \textbf{ALA-2}, la distribución
\(\phi/\psi\) pasa de un clúster compacto centrado en el cuadrante
\(\phi<0^\circ,\ \psi>0^\circ\) a \SI{298}{K} a una región igualmente estrecha pero
desplazada hacia \(\psi<0^\circ\) a \SI{400}{K}. Aunque la nube se alarga y la
densidad se reparte con mayor amplitud, el residuo sigue ocupando un único mínimo
conformacional, mostrando solo un aumento moderado de la flexibilidad torsional.

En \textbf{ARG-3}, a \SI{298}{K} la población se concentra de forma muy definida en la
región \(\phi\approx -120^\circ,\ \psi\approx 130{-}170^\circ\), formando un clúster alto y
bien delimitado. A \SI{400}{K}, el residuo conserva una única región poblada, pero la nube
se expande en ambas direcciones, con una dispersión mayor en \(\psi\) y contornos más
difusos ($\psi\approx -20$--$15^\circ$). Aun así, no se observan cambios de cuenca torsional: el residuo central
mantiene la misma familia conformacional, únicamente con fluctuaciones más amplias.

Por el contrario, \textbf{ASN-4} muestra la mayor sensibilidad térmica. A \SI{298}{K}, el mapa
presenta una región esencialmente continua en el cuadrante
\(\phi<0^\circ,\ \psi<0^\circ\). Sin embargo, a \SI{400}{K} la distribución se fragmenta en
dos mínimos bien diferenciados: uno alrededor de \(\psi\approx -60^\circ\) y otro cercano a
\(\psi\approx +10^\circ\), indicando un acceso más frecuente a configuraciones alternativas.
Este comportamiento bimodal refleja que el extremo C-terminal adquiere mayor
flexibilidad y explora distintas regiones permisibles del espacio \(\phi/\psi\) al aumentar
la temperatura.

Los mapas muestran que \textbf{ALA-2} y \textbf{ARG-3} mantienen una conformación
principal estable a ambas temperaturas, aunque más dispersa a \SI{400}{K}, mientras que
\textbf{ASN-4} experimenta una transición clara hacia una dinámica multimodal, siendo el
residuo más sensible al incremento térmico.

% ==================================================
\subsection{Puentes de hidrógeno}

\noindent\textbf{Nota}. En la ejecución de gmx hbond se seleccionaron como grupos donador y aceptor las opciones correspondientes a 'Protein' (número 1 y otra vez 1).

\begin{lstlisting}[style=bashstyle]
# Cuenta puentes de hidrógeno dentro del péptido
# -num: salida del número de H-bonds por frame
gmx hbond -f arn_run_298.trr -s arn_run_298.tpr -num hbond_298.xvg -xvg none

gmx hbond -f arn_run_400.trr -s arn_run_400.tpr -num hbond_400.xvg -xvg none
\end{lstlisting}

\subsubsection*{Análisis de los puentes de hidrógeno}

Durante el análisis interactivo, GROMACS identificó \textbf{8 donadores} y \textbf{13 aceptores} en el péptido. 
Aunque ello implica \textbf{104 combinaciones donador–aceptor} posibles en sentido teórico, 
el programa sólo considera como \emph{“pares posibles”} aquellos que cumplen los criterios 
geométricos iniciales (distancia \(\leq 0.35\,\mathrm{nm}\) y accesibilidad espacial en el grid de búsqueda). 
Por este motivo, GROMACS informa de un máximo de \textbf{52 pares geométricamente evaluables}, 
número inferior al total combinatorio. Aun así, el número real de enlaces de hidrógeno intramoleculares detectados 
es muy bajo, como se aprecia inspeccionando los ficheros generados con \texttt{less}.

\paragraph*{Inspección directa mediante \texttt{less}.}

A \textbf{298 K}, el archivo \texttt{hbond\_num\_298.xvg} muestra que la columna de H-bonds internos permanece
igual a cero en todos los frames:

\begin{verbatim}
   time     H-bonds   pairs
     0.000      0        18
     0.001      0        18
     0.007      0        19
     0.015      0        18
     ...
     1.995      0        21
     1.999      0        21
     2.000      0        21
\end{verbatim}

A \textbf{400 K}, el archivo \texttt{hbond\_num\_400.xvg} muestra la presencia de contactos transitorios,
con valores de 1–2 H-bonds en la mayoría de frames:

\begin{verbatim}
   time     H-bonds   pairs
     0.000      1        19
     0.006      1        20
     0.014      1        20
     0.017      1        19
     ...
     1.982      1        22
     1.991      0        23
     1.999      0        22
     2.000      0        21
\end{verbatim}

\paragraph*{Resumen cuantitativo.}

Con base en los valores extraídos por GROMACS:

\begin{itemize}
    \item A \textbf{298 K}, el número medio de enlaces de hidrógeno intramoleculares es \textbf{0.00}. 
    Ningún frame presenta enlaces internos estables, lo que es esperable en un tripéptido pequeño y
    completamente expuesto al solvente.

    \item A \textbf{400 K}, el número medio aumenta a \textbf{1.15 H-bonds por frame}. 
    Estos contactos son \textbf{transitorios} y reflejan configuraciones momentáneas favorecidas por la
    mayor flexibilidad térmica, sin llegar a formar una red interna estable.
\end{itemize}

\begin{figure}[H]
\centering
\includegraphics[width=0.90\textwidth]{img/figura_python_hbonds_298K.png}
\caption{Puentes de hidrógeno a 298 K.}
\end{figure}

\begin{figure}[H]
\centering
\includegraphics[width=0.90\textwidth]{img/figura_python_hbonds_400K.png}
\caption{Puentes de hidrógeno a 400 K.}
\end{figure}

\paragraph*{Análisis de los enlaces de hidrógeno intramoleculares}

Las figuras anteriores muestran la evolución temporal del número de enlaces de hidrógeno
\textbf{intramoleculares} del tripéptido ARN calculados con \texttt{gmx hbond -num}, donde la segunda columna
del archivo \texttt{.xvg} corresponde al número real de puentes de hidrógeno detectados por frame.

A \SI{298}{K}, la señal permanece constante en \textbf{cero} durante toda la trayectoria de \SI{2}{ps}. Esto confirma
que el péptido no forma enlaces de hidrógeno internos estables en condiciones térmicas moderadas, lo que es
esperable en un sistema tan corto y completamente expuesto al solvente.

A \SI{400}{K}, aparecen \textbf{contactos transitorios} con valores entre 0 y 2 enlaces por frame, con una media
de aproximadamente 1.15. Estos eventos corresponden a alineamientos momentáneos entre grupos N–H y
carbonilos del propio péptido favorecidos por la mayor flexibilidad térmica. No obstante, los contactos carecen
de persistencia temporal, y en ningún caso se configura una red interna estable.

Los resultados muestran que el tripéptido ARN no mantiene puentes de hidrógeno internos a
\SI{298}{K}, mientras que a \SI{400}{K} sólo se observan interacciones efímeras inducidas por la mayor agitación
térmica, sin comprometer la ausencia de estructura secundaria y el carácter altamente flexible del sistema.

% ==================================================
\subsection{Energía y temperatura}

\noindent\textbf{Nota}. gmx energy requiere elegir las magnitudes por número en el menú interactivo. En este trabajo, según la numeración mostrada por GROMACS, se seleccionaron:
14 -> Kinetic En. (energía cinética) ; 
15 -> Total Energy (energía total)
Para extraer la temperatura, en el menú de gmx energy se seleccionó:
17 -> Temperature (temperatura del sistema). 
La extracción se confirmó pulsando Enter tras introducir los números

\begin{lstlisting}[style=bashstyle]
# Extrae energía cinética, total y otros términos del fichero de energías (.edr)
# -f: fichero de energías generado en el RUN (registra magnitudes termodinámicas por paso)
# -s: .tpr con la definición del sistema (grupos, masas, g.d.l.); ayuda a etiquetar correctamente
# -o: salida en formato XVG (lista de tiempo vs magnitud) para graficar/análisis posterior
gmx energy -f arn_run_298.edr -s arn_run_298.tpr -o energia_298.xvg
gmx energy -f arn_run_400.edr -s arn_run_400.tpr -o energia_400.xvg

# Extrae la temperatura media y su serie temporal desde el .edr (tal y como la reporta el termostato)
# -o temp_XXX.xvg: serie de T(t); la media (y RMSD) aparece en la cabecera/estadísticos del comando
gmx energy -f arn_run_298.edr -s arn_run_298.tpr -o temp_edr_298.xvg
gmx energy -f arn_run_400.edr -s arn_run_400.tpr -o temp_edr_400.xvg
\end{lstlisting}

\subsubsection*{Análisis de la energía cinética, energía total y temperatura}

\begin{figure}[H]
\centering
\includegraphics[width=0.90\textwidth]{img/figura_python_energias_298K.png}
\caption{Energía cinética y energía total a 298 K.}
\end{figure}

\begin{figure}[H]
\centering
\includegraphics[width=0.90\textwidth]{img/figura_python_energias_400K.png}
\caption{Energía cinética y energía total a 400 K.}
\end{figure}

\paragraph*{Comparación de los perfiles energéticos a 298 K y 400 K.}

Las Figuras anteriores muestran la evolución temporal de la energía cinética y de
la energía total durante los 2 ps de producción NVT para ambas temperaturas. A
\SI{298}{K}, la energía cinética se mantiene estable alrededor de 
$\sim 9.7\times10^{3}$~kJ/mol, con fluctuaciones pequeñas y regulares.
La energía total se mantiene prácticamente plana en torno a
$-2.64\times10^{4}$~kJ/mol, lo que confirma la estabilidad del integrador y del
termostato en este régimen térmico.

A \SI{400}{K}, la energía cinética aumenta de forma proporcional a la
temperatura, estabilizándose en torno a $\sim 1.3\times10^{4}$~kJ/mol. La
amplitud de las fluctuaciones es mayor que a 298~K, coherente con la mayor
agitación térmica del sistema. La energía total presenta valores menos negativos
($\sim -1.82\times10^{4}$~kJ/mol), reflejando el incremento de la energía
interna. A pesar del aumento térmico, ambas señales se mantienen estables, 
lo que certifica que el termostato \texttt{v-rescale} preserva
adecuadamente el ensamble NVT en las dos condiciones.

Las gráficas confirman que el incremento de temperatura produce el
comportamiento esperado según el teorema de equipartición, sin comprometer la
estabilidad dinámica del sistema.

\begin{figure}[H]
\centering
\includegraphics[width=0.90\textwidth]{img/figura_python_temperatura_298K.png}
\caption{Evolución temporal de la temperatura a \SI{298}{K}}
\end{figure}

\begin{figure}[H]
\centering
\includegraphics[width=0.90\textwidth]{img/figura_python_temperatura_400K.png}
\caption{Evolución temporal de la temperatura a \SI{400}{K}}
\end{figure}

\paragraph*{Comparación de la estabilidad térmica a 298 K y 400 K.}

Las Figuras anteriores muestran la evolución temporal de la temperatura
instantánea durante los \SI{2}{ps} de producción NVT, permitiendo evaluar 
la respuesta del termostato \texttt{v-rescale} y la estabilidad térmica
del sistema. A \SI{298}{K}, la señal oscila de manera acotada entre
aproximadamente 285 y 310~K, con fluctuaciones típicas de $\pm 5{-}7$~K. El
trazado presenta pequeñas variaciones por intervalos de tiempo: un régimen
inicial ligeramente más disperso (0.0–0.3~ps), seguido por una zona más
estabilizada (0.3–0.75~ps), un periodo de ligera expansión térmica
(0.75–1.10~ps) y un tramo final donde la temperatura vuelve a concentrarse en
torno a 295–300~K. En todos los casos, la media se mantiene alineada con el
valor objetivo, sin deriva apreciable.

A \SI{400}{K}, la amplitud de las fluctuaciones aumenta de forma notable,
situándose habitualmente en el intervalo 380–420~K, con variaciones típicas de
$\pm 12{-}15$~K. Este aumento de varianza es coherente con la mayor energía
cinética disponible y con la dinámica interna más activa del tripéptido. La
serie temporal muestra regiones con oscilaciones especialmente intensas
(1–1.40~ps), que reflejan configuraciones más energéticas, y
tramos donde la señal se estrecha. Sin embargo,
el valor medio permanece centrado en los \SI{400}{K} durante toda la
trayectoria.

Ambas gráficas confirman que el termostato \texttt{v-rescale}
mantiene la temperatura objetivo con alta fidelidad en los dos regímenes, con un
incremento de la dispersión proporcional a la temperatura según la teoría de
equipartición. 

\subsubsection*{Resumen cuantitativo de producción (NVT, 2 ps)}

\begin{table}[H]
\centering
\begin{tabular}{lrrrr}
\toprule
\textbf{Magnitud} & \textbf{298 K (media)} & \textbf{Err.Est.} & \textbf{RMSD} & \textbf{Tot-Drift} \\
\midrule
Kinetic Energy (kJ/mol) & 9735.05 & 39  & 155.649 & -148.396 \\
Total Energy (kJ/mol)   & -26443.6 & 110 & 263.757 & -642.346 \\
Temperature (K)         & 296.008  & 1.2 & 4.733   & -4.512 \\
\bottomrule
\end{tabular}
\caption{Energía y temperatura en producción NVT a 298 K (2 ps).}
\end{table}

\begin{table}[H]
\centering
\begin{tabular}{lrrrr}
\toprule
\textbf{Magnitud} & \textbf{400 K (media)} & \textbf{Err.Est.} & \textbf{RMSD} & \textbf{Tot-Drift} \\
\midrule
Kinetic Energy (kJ/mol) & 13077.4 & 45  & 192.186 & 104.161 \\
Total Energy (kJ/mol)   & -18174.6 & 75 & 228.932 & 83.218 \\
Temperature (K)         & 397.636  & 1.4 & 5.844  & 3.167 \\
\bottomrule
\end{tabular}
\caption{Energía y temperatura en producción NVT a 400 K (2 ps).}
\end{table}

\subsubsection{Validación teórica de la energía cinética (Equipartición)}

Para comprobar que las simulaciones en ensamble NVT son físicamente correctas,
se comparó la energía cinética media reportada por GROMACS con el valor teórico
esperado según el teorema de equipartición:

\begin{equation}
E_k = \frac{3}{2} \, N \, k_B T = \frac{3}{2} N R T,
\end{equation}

Donde $N$ es el número de átomos del sistema ($N=2638$), $T$ la temperatura media
obtenida en la simulación y $R = 8.314\times 10^{-3}$ kJ mol$^{-1}$ K$^{-1}$ la
constante de los gases.

\paragraph{Simulación a 298 K.}
GROMACS proporcionó una energía cinética media de $E_k^{\mathrm{GMX}} = 9735.05$
kJ/mol y una temperatura media de $T = 296.008$ K. El valor teórico es:

\begin{equation}
E_k^{\mathrm{teo}} = \frac{3}{2} \, N R T
= 2638 \times \frac{3}{2} \times (8.314\times 10^{-3}) \times 296.008
= 9735.0\ \text{kJ/mol}.
\end{equation}

La diferencia relativa es menor del $10^{-4}$\%, lo cual demuestra una excelente
coherencia entre el termostato y el teorema de equipartición.

\paragraph{Simulación a 400 K.}
GROMACS reportó $E_k^{\mathrm{GMX}} = 13077.4$ kJ/mol y una temperatura media de
$T = 397.636$ K. El valor teórico correspondiente es:

\begin{equation}
E_k^{\mathrm{teo}} = \frac{3}{2} \, N R T
= 2638 \times \frac{3}{2} \times (8.314\times 10^{-3}) \times 397.636
= 13159.4\ \text{kJ/mol}.
\end{equation}

La desviación es $\sim 0.6\%$, perfectamente aceptable para una trayectoria corta
y totalmente compatible con el comportamiento esperado del ensamble NVT.

\paragraph{Conclusión.}
En ambos casos, la energía cinética simulada concuerda con la predicción teórica.

% ==================================================
\subsection{Velocidades atómicas}

\noindent\textbf{Átomos seleccionados.}
Se analizaron cinco átomos contiguos del esqueleto peptídico:
N(ALA-2)\,[7], CA(ALA-2)\,[9], C(ALA-2)\,[15], N(ARG-3)\,[17] y CA(ARG-3)\,[19].

\noindent\textbf{Extracción con GROMACS.}
Se generó un fichero de índice con esos cinco átomos y se extrajeron sus
velocidades (módulo y componentes) a partir de la trayectoria:

\begin{lstlisting}[style=bashstyle]
cat > cinco_atoms.ndx << 'EOF'
[ cinco ]
7 9 15 17 19
EOF
# Extrae velocidades atómicas desde la trayectoria
# -ov genera un archivo con |v| y componentes vx, vy, vz
# 298 K (solo esos 5 átomos)
gmx traj -f arn_run_298.trr -s arn_run_298.tpr -n cinco_atoms.ndx -ov vel5_298.xvg

# 400 K (solo esos 5 átomos)
gmx traj -f arn_run_400.trr -s arn_run_400.tpr -n cinco_atoms.ndx -ov vel5_400.xvg
\end{lstlisting}

\subsubsection*{Análisis de las velocidades atómicas}

\paragraph*{Salidas y formato de datos.}
Los comandos anteriores generan los ficheros:
\texttt{vel5\_298.xvg} y \texttt{vel5\_400.xvg}, en formato \emph{XVG} (GROMACS/Grace).
Cada archivo contiene:
(i) un encabezado con metadatos y leyendas (\texttt{@ s\# legend "..."}),
(ii) columnas numéricas con las series temporales.
El contenido de las columnas es:
\[
\text{col}_0 = t~(\mathrm{ps}),\quad
\big(\,v_x^{(a_1)},\, v_y^{(a_1)},\, v_z^{(a_1)}\big),\;
\ldots,\;
\big(\,v_x^{(a_5)},\, v_y^{(a_5)},\, v_z^{(a_5)}\big),
\]
con \(\mathrm{nm\,ps^{-1}}\) como unidades de velocidad. Los identificadores de átomo en las leyendas (\texttt{"atom 7 X"}, \texttt{"atom 9 Y"}, etc.) confirman el orden y permiten recuperar el ID de cada átomo analizado. En el postproceso, el módulo \(|\mathbf{v}|\) de cada átomo se calcula como
\(|\mathbf{v}|=\sqrt{v_x^2+v_y^2+v_z^2}\) para construir las figuras.

\paragraph*{Nota.}Las velocidades atómicas proporcionan información directa sobre la agitación térmica de los átomos y permiten verificar el efecto del aumento de temperatura sobre la dinámica interna del tripéptido. Se analizaron cinco átomos consecutivos del \emph{backbone}: N(ALA–2) [7], CA(ALA–2) [9], C(ALA–2) [15], N(ARG–3) [17] y CA(ARG–3) [19]. Para cada átomo se extrajeron las componentes $v_x, v_y, v_z$ y se calculó el módulo $|\mathbf{v}|$ (mediante la expresión $|\mathbf{v}|=\sqrt{v_x^2+v_y^2+v_z^2}$) en unidades de \(\mathrm{nm\,ps^{-1}}\).

\begin{figure}[H]
  \centering
  \includegraphics[width=0.94\textwidth]{img/figura_python_vel_5atomos_298K.png}
  \caption{\textbf{Velocidades atómicas a \SI{298}{K}}.
  Evolución temporal del módulo de la velocidad $|\mathbf{v}|$ para los cinco átomos del \emph{backbone}.
  Estadísticos: 
  N(ALA–2) [7] $\langle |v| \rangle = \SI{0.654}{nm.ps^{-1}}$, $\sigma = \SI{0.266}{nm.ps^{-1}}$;
  CA(ALA–2) [9] $\langle |v| \rangle = \SI{0.619}{nm.ps^{-1}}$, $\sigma = \SI{0.254}{nm.ps^{-1}}$;
  C(ALA–2) [15] $\langle |v| \rangle = \SI{0.616}{nm.ps^{-1}}$, $\sigma = \SI{0.264}{nm.ps^{-1}}$;
  N(ARG–3) [17] $\langle |v| \rangle = \SI{0.626}{nm.ps^{-1}}$, $\sigma = \SI{0.262}{nm.ps^{-1}}$;
  CA(ARG–3) [19] $\langle |v| \rangle = \SI{0.681}{nm.ps^{-1}}$, $\sigma = \SI{0.306}{nm.ps^{-1}}$.
  }
  \label{fig:vel5-298}
\end{figure}

\begin{figure}[H]
  \centering
  \includegraphics[width=\textwidth]{img/fig_vel_5atomos_298K_panel.png}
  \caption{\textbf{Panel 5×1 del módulo de la velocidad a \SI{298}{K}}.
  Cada panel muestra el comportamiento de un átomo específico del \emph{backbone},
  facilitando la comparación visual de sus fluctuaciones individuales.}
  \label{fig:vel5-298-panel}
\end{figure}

\begin{figure}[H]
  \centering
  \includegraphics[width=0.94\textwidth]{img/figura_python_vel_5atomos_400K.png}
  \caption{\textbf{Velocidades atómicas a \SI{400}{K}}.
  Se observa un aumento claro en la amplitud de las fluctuaciones respecto a \SI{298}{K}, coherente con el incremento térmico.
  Estadísticos:
  N(ALA–2) [7] $\langle |v| \rangle = \SI{0.996}{nm.ps^{-1}}$, $\sigma = \SI{0.426}{nm.ps^{-1}}$;
  CA(ALA–2) [9] $\langle |v| \rangle = \SI{0.960}{nm.ps^{-1}}$, $\sigma = \SI{0.391}{nm.ps^{-1}}$;
  C(ALA–2) [15] $\langle |v| \rangle = \SI{0.735}{nm.ps^{-1}}$, $\sigma = \SI{0.302}{nm.ps^{-1}}$;
  N(ARG–3) [17] $\langle |v| \rangle = \SI{0.777}{nm.ps^{-1}}$, $\sigma = \SI{0.324}{nm.ps^{-1}}$;
  CA(ARG–3) [19] $\langle |v| \rangle = \SI{0.734}{nm.ps^{-1}}$, $\sigma = \SI{0.306}{nm.ps^{-1}}$.
  }
  \label{fig:vel5-400}
\end{figure}

\begin{figure}[H]
  \centering
  \includegraphics[width=\textwidth]{img/fig_vel_5atomos_400K_panel.png}
  \caption{\textbf{Panel 5×1 del módulo de la velocidad a \SI{400}{K}}.
  El aumento de amplitud en todos los átomos refleja la mayor energía cinética media
  y la dispersión esperable en un régimen térmico elevado.}
  \label{fig:vel5-400-panel}
\end{figure}

\paragraph*{Comparación entre 298 K y 400 K.}

Las Figuras~\ref{fig:vel5-298}--\ref{fig:vel5-400-panel} muestran un 
incremento claro y sistemático de la agitación térmica al pasar de \SI{298}{K} a 
\SI{400}{K}. En ambos regímenes las señales presentan oscilaciones rápidas y 
centradas en cero, sin deriva de momento, lo que confirma la correcta actuación 
del termostato. No obstante, la amplitud de las fluctuaciones y la dispersión 
estadística aumentan de manera apreciable con la temperatura.

A \SI{298}{K}, los cinco átomos presentan módulos de velocidad relativamente 
compactos, con valores medios en el intervalo 
\(0.62{-}0.68~\mathrm{nm\,ps^{-1}}\). El átomo más móvil es CA(ARG–3) [19], con 
\(\langle|v|\rangle=\SI{0.681}{nm.ps^{-1}}\), mientras que el resto oscila en torno 
a \(0.62{-}0.65\). Las desviaciones estándar, próximas a 
\(0.25{-}0.30~\mathrm{nm\,ps^{-1}}\), indican fluctuaciones moderadas y bien 
contenidas, coherentes con un régimen térmico suave.

A \SI{400}{K}, todos los átomos incrementan su velocidad media y su dispersión. 
Los valores medios oscilan ahora entre 
\(0.73{-}1.00~\mathrm{nm\,ps^{-1}}\). El átomo más móvil pasa a ser N(ALA–2) [7], 
con \(\langle|v|\rangle=\SI{0.996}{nm.ps^{-1}}\), y presenta además la mayor 
desviación estándar del conjunto (\(\sigma=\SI{0.426}{nm.ps^{-1}}\)), reflejando 
una agitación térmica más intensa. El resto de átomos sigue la misma tendencia, 
con incrementos relativos del \(20{-}50\%\) respecto a \SI{298}{K}. Este aumento 
se aprecia con especial claridad en las figuras de tipo panel, donde las 
señales muestran transiciones más amplias y picos más frecuentes, sin pérdida 
de estabilidad numérica ni comportamientos no físicos.

El aumento de temperatura induce un incremento global y homogéneo de las
velocidades atómicas. A \SI{298}{K}, las diferencias entre átomos son muy
pequeñas y no aparece un patrón claro, con valores medios muy próximos entre
sí. En cambio, a \SI{400}{K} sí se observa el comportamiento esperado: los
átomos más ligeros y expuestos del backbone, especialmente los átomos N,
se convierten en los más rápidos, mientras que los carbonos centrales
mantienen valores más amortiguados.

% ==================================================
\subsection{Distribución de temperatura}

\noindent\textbf{Nota Temperatura.} La temperatura instantánea se extrajo a partir de la
trayectoria completa (\texttt{.trr}) utilizando el grupo \texttt{System}. La producción
tiene una duración total de \(2~\mathrm{ps}\) con un intervalo temporal de
\(1~\mathrm{fs}\), lo que genera exactamente \(2001\) frames de temperatura. Esta
frecuencia de muestreo fija la anchura estadística esperada de la distribución,
ya que los valores consecutivos están fuertemente correlacionados en una
trayectoria tan corta.

\noindent\textbf{Nota menú}. En el menú se seleccionó \texttt{System} (opción 0), ya que la temperatura debe calcularse sobre todos los átomos del sistema.

\begin{lstlisting}[style=bashstyle]
# Extrae la temperatura instantánea por frame a partir de la trayectoria de 298 K.
# -f arn_run_298.trr : trayectoria de producción (formato .trr, incluye velocidades)
# -s arn_run_298.tpr : archivo de configuración/topología con masas y g.d.l. (necesario para T)
# -ot temp_series_298.xvg : salida con T(t) (temperatura vs tiempo) en formato XVG

gmx traj -f arn_run_298.trr -s arn_run_298.tpr -ot temp_298.xvg

# Lo mismo para 400 K:
gmx traj -f arn_run_400.trr -s arn_run_400.tpr -ot temp_400.xvg
\end{lstlisting}

\subsubsection*{Análisis de la distribución de temperatura}

\begin{figure}[H]
\centering
\includegraphics[width=0.90\textwidth]{img/figura_python_histograma_298K.png}
\caption{Distribución estadística de T instantánea a 298 K.}
\end{figure}

\begin{figure}[H]
\centering
\includegraphics[width=0.90\textwidth]{img/figura_python_histograma_400K.png}
\caption{Distribución estadística de T instantánea a 400 K.}
\end{figure}

\paragraph*{Análisis de las distribuciones de temperatura}

Las distribuciones de temperatura obtenidas a \SI{298}{K} y \SI{400}{K} presentan un
comportamiento claramente gaussiano, en línea con el esperado para un ensamble
NVT controlado con el termostato \texttt{v-rescale}. 

A \SI{298}{K}, la distribución es estrecha y centrada en \(\sim 297\,\mathrm{K}\), lo que
indica una fluctuación térmica moderada y una correcta estabilización del sistema.
A \SI{400}{K}, la anchura de la distribución aumenta de forma apreciable, reflejando
la mayor agitación térmica; aun así, el máximo permanece en \(\sim 396\,\mathrm{K}\),
confirmando la ausencia de derivas y el mantenimiento estable de la temperatura
objetivo.

% =============================
% 6. Conclusiones
% =============================
\section{Conclusiones}
El análisis comparado de las simulaciones a \SI{298}{K} y \SI{400}{K} muestra que el
aumento de temperatura produce un incremento en la movilidad interna del
tripéptido ARN. A \SI{400}{K}, el radio de giro presenta fluctuaciones más amplias
y los diedros del backbone exploran regiones conformacionales más extensas,
reflejando una mayor flexibilidad torsional. Los puentes de hidrógeno
intramoleculares que aparecen a esta temperatura son transitorios y de muy baja
persistencia, sin llegar a configurar una red estable.

Las distancias covalentes y los ángulos internos del esqueleto peptídico se
mantienen parecidas en ambas condiciones térmicas, lo que
evidencia la rigidez geométrica del backbone. Los mapas de Ramachandran muestran
una dispersión mayor a \SI{400}{K}, especialmente en el residuo ASN-4, aunque las
conformaciones siguen confinadas en regiones permitidas del espacio
\(\phi/\psi\), sin acceso a estados desfavorables.

La energía cinética media concuerda con los valores previstos por el teorema de
equipartición en ambos regímenes, validando el funcionamiento del termostato
\texttt{v-rescale} y la estabilidad del ensamble NVT. La evolución de la energía
total y de la temperatura instantánea confirma trayectorias sin deriva y una
integración numérica estable.

En conjunto, el tripéptido ARN conserva su arquitectura global tanto a
\SI{298}{K} como a \SI{400}{K}, aunque a mayor temperatura accede a un espacio
conformacional más amplio y flexible. Estos resultados muestran el impacto de la
temperatura sobre la dinámica interna de pequeños péptidos en disolución.
\end{document}

